\documentclass{article}

\input{header}

\title{CSCI 2590: Assignment 4}

\author{Full Name \\ Net ID \\ Your Github Link}

\date{}


\colmfinalcopy
\begin{document}
\maketitle
% \section*{Part I. Q1} No written element, submit \texttt{out\_original.txt}  to autograder.
\section*{Q0. 1.}
Please provide a link to your github repository, which contains the code for both Part I and Part II. \textcolor{gray}{TODO}
\section*{Q2. 1.}
Describe your transformation of dataset.
We design a \textbf{label-preserving noisy text transformation} that simulates realistic user writing errors in movie reviews. The transformation is applied to each review independently with two operations:

\begin{itemize}
  \item \textbf{Synonym substitution:} For each eligible content word, with probability $p_s=0.12$, replace it with a WordNet synonym (if available). We preserve casing and skip replacements that change token type (e.g., punctuation/numbers).
  \item \textbf{Keyboard-typo injection:} For each eligible word, with probability $p_t=0.08$, replace one character by a nearby QWERTY key (e.g., \texttt{e}$\rightarrow$\texttt{r}, \texttt{i}$\rightarrow$\texttt{o}).
\end{itemize}

To reduce label flips, we do \textbf{not} modify negation cues and sentiment-critical markers such as \textit{not, never, no, n't}. We also skip very short words (length $<3$), stopwords, and special tokens. This keeps transformed text grammatical enough for humans while still introducing distribution shift for the model.

This transformation is \textbf{reasonable} because real users commonly produce exactly these variations at test time: minor typos and lexical paraphrases. At the same time, it is non-trivial for the classifier, since token-level perturbations change surface forms and can degrade robustness without changing underlying sentiment.

Examples:
\begin{itemize}
  \item Original: ``This movie is absolutely amazing and very entertaining.''
  \item Transformed: ``This film is absolutely amazinh and very entertaining.''
  \item Original: ``I did not like this boring plot at all.''
  \item Transformed: ``I did not like this dull plot at all.''
\end{itemize}
% \section*{Part I. Q2. 2. No written element, submit \texttt{out\_transformed.txt} to autograder. }
\section*{Q3. 1}
\textbf{Report \& Analysis}
    \begin{itemize}
    \item Reported accuracies: original baseline model on original test set = 0.93116; baseline model on transformed test set = 0.90892; augmented model on original test set = 0.92968; augmented model on transformed test set = 0.91504.
    \item The transformed-set performance improved after augmentation (0.90892 $\rightarrow$ 0.91504, +0.00612 absolute). On the original set, augmentation caused a small decrease (0.93116 $\rightarrow$ 0.92968, -0.00148 absolute).
    \item Intuitively, augmentation exposes the model to perturbations similar to the evaluation-time distribution shift, so the classifier becomes less sensitive to lexical/surface variations. This improves robustness on transformed inputs, but slightly trades off specialization to the clean original distribution.
    \item One limitation is coverage mismatch: the augmented examples reflect only one handcrafted perturbation family. If real OOD inputs differ (e.g., different typo patterns, syntax changes, or paraphrases), gains may not transfer well.
    \end{itemize}
\section*{Part II. Q4}
% 
% \section{Data Statistics and Processing (8pt)}


\begin{table}[h!]
\centering
\begin{tabular}{lcc}
\toprule
Statistics Name & Train & Dev \\
\midrule
Number of examples & 4225 & 466 \\
Mean sentence length & 17.10 & 17.07 \\
Mean SQL query length & 216.37 & 210.05  \\
Vocabulary size (natural language)& 791 & 465  \\
Vocabulary size (SQL)& 555 & 395  \\
\bottomrule
\end{tabular}
\caption{Data statistics before any pre-processing. \textcolor{gray}{You need to at least provide the statistics listed above, and can add new entries.}}
\label{tab:data_stats_before}
\end{table}

\begin{table}[h!]
\centering
\begin{tabular}{lcc}
\toprule
\textbf{T5 fine-tuned model} & Train & Dev \\
\midrule
Mean sentence length & 22.10 & 22.07 \\
Mean SQL query length & 217.37 & 211.05 \\
Vocabulary size (natural language) & 795 & 469 \\
Vocabulary size (SQL) & 556 & 396 \\
\midrule
\bottomrule
\end{tabular}
\caption{Data statistics after pre-processing. \textcolor{gray}{You need to at least provide the statistics listed in \autoref{tab:data_stats_before} (except for the number of lines), and can add new entries.}}
\label{tab:data_stats_after}
\end{table}



\newpage




\section*{Q5}\label{sec:t5}


\begin{table}[h!]
\centering
\begin{tabular}{p{3.5cm}p{10cm}}
\toprule
Design choice & Description \\
\midrule
Data processing & \textcolor{gray}{Describe the data processing steps you undertook, if any.} \\
Tokenization & \textcolor{gray}{Describe how you did the tokenization for the inputs to the encoder and decoder. If you use anything else than the default T5 tokenizer, specify what you use, and why you choose to use it.} \\
Architecture & \textcolor{gray}{Describe the components in the T5 architecture that you chose to fine-tune. Did you fine-tune the entire model, specific layers?} \\
Hyperparameters & \textcolor{gray}{List the key hyperparameters that you used, including the learning rate, batch size, and stopping criteria.} \\
\bottomrule
\end{tabular}
\caption{Details of the best-performing T5 model configurations (fine-tuned)}
\label{tab:t5_results_ft}
\end{table}







\section*{Q6. }
\label{Q6}

\paragraph{Quantitative Results:} 
\begin{table}[h!]
\centering
\begin{tabular}{lcc}
  \toprule
  System & Query EM & F1 score\\
  \midrule
  \multicolumn{3}{l}{\textbf{Dev Results}} \\
  \midrule
  
  \multicolumn{3}{l}{T5 fine-tuned} \\
  Fully Trained model (from scratch) & XX.XX & XX.XX \\[5pt]
  % Variant1 & XX.XX & XX.XX \\
  % Variant2 & XX.XX & XX.XX \\
  % Variant3 & XX.XX & XX.XX \\
  
  \midrule
  \multicolumn{3}{l}{\textbf{Test Results}} \\
  \midrule
  T5 fine-tuning & XX.XX & XX.XX \\
Fully Trained model (from scratch) & XX.XX & XX.XX \\
  \bottomrule
\end{tabular}  
\caption{Development and test results. \textcolor{gray}{Use this table to report quantitative results for both dev and test results.}}
\label{tab:results}
\end{table}


\paragraph{Qualitative Error Analysis:} 

\begin{table}
  \centering
  \begin{tabular}{p{2cm}p{3cm}p{3cm}p{5cm}}
    \toprule
    \textbf{Error Type}& \textbf{Example Of Error} & \textbf{Error Description} & \textbf{Statistics} \\
    \midrule
    \textcolor{gray}{Error name}  & \textcolor{gray}{Snippet from datapoint examplifying error} & \textcolor{gray}{Describe the error in natural language} & \textcolor{gray}{Provide statistics in the form ``COUNT/TOTAL'' on the prevalence of the error. TOTAL is the number of relevant examples (e.g. number of queries, for query-level error), and COUNT is the number of examples that showed this error.}  \\
    
    \midrule
    &  &   & \\
    % \bottomrule
  \end{tabular}
  \label{tab:qualitative}
  \caption{Use this table for your qualitative analysis on the dev set.}\label{tab: qualitative}
\end{table}


\section*{Q7.}

Provide a link to a Google Drive that contains a model checkpoint used to generate outputs you have submitted. 
\textcolor{gray}{TODO}

\section*{Extra Credit: }
\textcolor{gray}{[Optional]}
If you are doing an extra credit assignment, please describe your system here, as well as provide a link to a Google Drive that contains a model checkpoint used to generate outputs you have submitted. Report the same metrics you reported in \ref{Q6}Q6



\section*{AI Usage}


\noindent\textbf{1. AI Tools Used} \\
(e.g., ChatGPT, Claude, Gemini) \\
\hrulefill



\noindent\textbf{2. How AI Was Used} \\
Describe how you used AI (e.g., debugging, explaining concepts, generating code). \\
\hrulefill


\noindent\textbf{3. Example Prompts / Interactions} \\
Provide 1--3 representative prompts or queries you used. \\
\hrulefill


\section*{Extra Credit: Vibe Jam user study}
\textcolor{gray}{[Optional]}
If you are doing Vibe Jam user study, please mention the following in gradescope in the Gradescope Assignment titled \textbf{``CS2590-Vibe Jam User Study"}
\begin{itemize}
\item Email
\item The user study completion code.
\item The screenshot of compensation checklist.
\end{itemize}






\end{document}